% !TEX program = xelatex
\documentclass[a4paper,10pt]{article}

\makeatletter
\let\@expl@@@mark@update@singlecol@structures@@\relax
\makeatother

\usepackage[a4paper, left=0.95cm, right=0.95cm, top=0.85cm, bottom=0.92cm, footskip=0.52cm]{geometry}
\usepackage{fontspec}
\usepackage{amsmath, amssymb}
\usepackage{longtable, array, colortbl}
\usepackage[table]{xcolor}
\usepackage{fancyhdr}
\usepackage{tikz}
\usetikzlibrary{calc, angles, quotes, arrows.meta, decorations.markings}

% Cross-platform font discovery
\makeatletter
\newcommand{\NGSFindFonts}{%
  \IfFileExists{../../../fonts/KhmerOSsiemreap.ttf}{%
    \def\NGSUserPath{../../../fonts/}\def\NGSSysPath{../../../fonts/}%
  }{%
    \IfFileExists{../../../../fonts/KhmerOSsiemreap.ttf}{%
      \def\NGSUserPath{../../../../fonts/}\def\NGSSysPath{../../../../fonts/}%
    }{%
      \IfFileExists{fonts/KhmerOSsiemreap.ttf}{%
        \def\NGSUserPath{fonts/}\def\NGSSysPath{fonts/}%
      }{%
        \def\NGSUserPath{C:/Users/user/AppData/Local/Microsoft/Windows/Fonts/}%
        \def\NGSSysPath{C:/Windows/Fonts/}%
      }%
    }%
  }%
}
\NGSFindFonts
\makeatother

\XeTeXlinebreaklocale "km"
\XeTeXlinebreakskip = 0pt plus 1pt minus 0.1pt

\setmainfont[
  Path = {\NGSUserPath},
  Extension = .ttf,
  UprightFont = KhmerOSsiemreap,
  BoldFont = KhmerOSsiemreap,
  BoldFeatures = {FakeBold=1.8},
  ItalicFont = KhmerOSsiemreap,
  ItalicFeatures = {FakeSlant=0.25}
]{KhmerOSsiemreap}

\newfontfamily\muol[Path={\NGSUserPath}]{KhmerOSmuollight.ttf}
\newfontfamily\Tacteing[Path={\NGSUserPath}]{TACTENG.TTF}
\newfontfamily\enfont[Path={\NGSSysPath}, Scale=1.0]{times.ttf}
\newfontfamily\enbold[Path={\NGSSysPath}, Scale=1.0]{timesbd.ttf}
\newfontfamily\enital[Path={\NGSSysPath}, Scale=1.0]{timesi.ttf}
\newfontfamily\calibrifont[Path={\NGSSysPath}, Scale=0.95]{calibri.ttf}
\newfontfamily\symfont[Path={\NGSSysPath}, Scale=0.95]{cour.ttf}

\DeclareTextFontCommand{\en}{\enfont}
\newcommand{\enbf}[1]{{\enbold #1}}
\newcommand{\enit}[1]{{\enital #1}}
\newcommand{\khmflourish}{{\Tacteing\fontsize{17pt}{17pt}\selectfont 3}}
\newcommand{\wcirc}{{\symfont\char"25CB}}

\definecolor{headgray}{HTML}{E6E6E6}
\definecolor{stepgray}{HTML}{D4D4D4}
\definecolor{fiveecolor}{HTML}{8B0000}

\pagestyle{fancy}
\fancyhf{}
\renewcommand{\headrulewidth}{0pt}
\renewcommand{\footrulewidth}{0pt}
\fancyfoot[R]{{\calibrifont\fontsize{9.5pt}{11pt}\selectfont\thepage}}

\setlength{\parindent}{0pt}
\setlength{\LTpre}{2pt}
\setlength{\LTpost}{2pt}
\setlength{\tabcolsep}{4.5pt}
\newcolumntype{L}[1]{>{\raggedright\arraybackslash}p{#1}}

\newcommand{\fiveE}[2]{{\color{fiveecolor}\enbf{#1} \enit{(#2)}}}
\newcommand{\timenote}[1]{{\color{fiveecolor}\textit{#1}}}
\newcommand{\stepheader}[1]{%
  \rowcolor{stepgray}%
  \multicolumn{3}{|c|}{\rule[-4pt]{0pt}{15pt}{\muol\fontsize{10pt}{14pt}\selectfont #1}} \\*%
  \hline
}
\newcommand{\cfig}[1]{\begin{center}\vspace{-2pt}#1\vspace{-3pt}\end{center}}

\newcommand{\tri}[1]{\triangle #1}
\newcommand{\unit}[1]{\,\text{\en{#1}}}
\newcommand{\cnum}[1]{\tikz[baseline=-0.58ex]\node[circle, fill=black, text=white, inner sep=0.85pt, minimum size=1.95ex, font=\enbold\scriptsize]{#1};}

\newcommand{\arc}[1]{\mbox{%
  \tikz[baseline=(arcX.base)]{%
    \node[inner sep=0pt, outer sep=0pt] (arcX) {$#1$};%
    \draw[black, line width=0.55pt] ($(arcX.north west)+(0.02,0.035)$) to[bend left=30] ($(arcX.north east)+(-0.02,0.035)$);%
  }%
}}

\newcommand{\tickone}[2]{%
  \coordinate (tkM) at ($(#1)!0.5!(#2)$);
  \draw[black, line width=0.7pt] ($(tkM)!0.11cm!90:(#2)$) -- ($(tkM)!0.11cm!-90:(#2)$);
}
\newcommand{\ticktwo}[2]{%
  \coordinate (tkM1) at ($(#1)!0.47!(#2)$);
  \coordinate (tkM2) at ($(#1)!0.53!(#2)$);
  \draw[black, line width=0.7pt] ($(tkM1)!0.11cm!90:(#2)$) -- ($(tkM1)!0.11cm!-90:(#2)$);
  \draw[black, line width=0.7pt] ($(tkM2)!0.11cm!90:(#2)$) -- ($(tkM2)!0.11cm!-90:(#2)$);
}

% ============================================================================
% ម៉ាក្រូរូបភាព TikZ ស-ខ្មៅ សមាមាត្រច្បាស់ (3.8 cm – 6.8 cm, អក្សរខ្មៅ 100%)
% ============================================================================

% ១. រូបបេសកកម្ម E1៖ ចង្អេរឫស្សី C(O, 25cm) និងការដាក់របារឫស្សីត្រង់វែងបំផុត
\newcommand{\figEngageCIRCThree}{%
\begin{tikzpicture}[scale=0.50, every node/.style={font=\footnotesize, text=black}]
  \coordinate (O) at (0,0);
  \def\R{2.8}
  \draw[black, line width=1.2pt] (O) circle (\R);
  \draw[black!45, line width=0.6pt] (O) circle (2.62);
  % អង្កត់ផ្ចិត AB (របារកាត់ផ្ចិត) និងអង្កត់ធ្នូផ្សេងៗ
  \draw[black, line width=1.15pt] (-\R,0) -- (\R,0);
  \draw[black, line width=0.9pt, densely dashed] (145:\R) -- (35:\R);
  \draw[black, line width=0.9pt, densely dashed] (-130:\R) -- (-30:\R);
  \fill[black] (O) circle (2.3pt) (-\R,0) circle (2.1pt) (\R,0) circle (2.1pt) (145:\R) circle (2pt) (35:\R) circle (2pt);
  \node[below=2pt] at (O) {ផ្ចិត $O$};
  \node[left=1.5pt] at (-\R,0) {$A$};
  \node[right=1.5pt] at (\R,0) {$B$};
  \node[above left=0.5pt] at (145:\R) {$C$};
  \node[above right=0.5pt] at (35:\R) {$F$};
  \node[above=1pt] at (0,0.02) {$AB = ?$};
  \node[above=1pt] at (0,1.62) {$CF$};
\end{tikzpicture}}

% ២. រូប E2៖ អង្កត់ផ្ចិតធំជាងគេ (AB > CD) និងអន្ទាក់មិនសមាមាត្រ (ធ្នូ CF = 2CD តែ CF < 2CD)
\newcommand{\figExploreCIRCThree}{%
\begin{tikzpicture}[scale=0.38, every node/.style={font=\scriptsize, text=black}]
  % ខាងឆ្វេង៖ អង្កត់ផ្ចិត AB > អង្កត់ធ្នូ CD (ទំ.២០១)
  \begin{scope}[shift={(0,0)}]
    \coordinate (O) at (3.0,2.8);
    \def\R{2.35}
    \coordinate (A) at ($(O)+(-\R,0)$);
    \coordinate (B) at ($(O)+(\R,0)$);
    \coordinate (C) at ($(O)+(135:\R)$);
    \coordinate (D) at ($(O)+(30:\R)$);
    \draw[black, line width=0.95pt] (O) circle (\R);
    \draw[black, line width=1.1pt] (A)--(B) (C)--(D);
    \draw[black, line width=0.8pt, densely dashed] (O)--(C) (O)--(D);
    \tickone{O}{C}\tickone{O}{D}
    \fill[black] (O) circle (2.2pt) (A) circle (2pt) (B) circle (2pt) (C) circle (2pt) (D) circle (2pt);
    \node[below=1.5pt] at (O) {$O$};
    \node[left=1pt] at (A) {$A$};
    \node[right=1pt] at (B) {$B$};
    \node[above left=0.5pt] at (C) {$C$};
    \node[above right=0.5pt] at (D) {$D$};
    \node[align=center] at (3.0,-1.65) {គំរូទី១៖ $OC+OD > CD$\\$\Rightarrow AB = 2r > CD$};
  \end{scope}
  % ខាងស្ដាំ៖ អន្ទាក់មិនសមាមាត្រ 60° និង 120°
  \begin{scope}[shift={(9.3,0)}]
    \coordinate (O2) at (3.0,2.8);
    \def\R{2.35}
    \coordinate (C2) at ($(O2)+(150:\R)$);
    \coordinate (D2) at ($(O2)+(90:\R)$);
    \coordinate (F2) at ($(O2)+(30:\R)$);
    \draw[black, line width=0.95pt] (O2) circle (\R);
    \draw[black, line width=1.05pt] (C2)--(D2)--(F2) (C2)--(F2);
    \draw[black, line width=0.75pt, densely dashed] (O2)--(C2) (O2)--(D2) (O2)--(F2);
    \tickone{C2}{D2}\tickone{D2}{F2}
    \fill[black] (O2) circle (2.2pt) (C2) circle (2pt) (D2) circle (2pt) (F2) circle (2pt);
    \node[below=1.5pt] at (O2) {$O$};
    \node[above left=0.5pt] at (C2) {$C$};
    \node[above=1pt] at (D2) {$D$};
    \node[above right=0.5pt] at (F2) {$F$};
    \node[align=center] at (3.0,-1.65) {អន្ទាក់៖ $\arc{CF} = 2\arc{CD}$\\តែ $CF < CD+DF = 2CD$};
  \end{scope}
\end{tikzpicture}}

% ៣. រូប E4៖ លំហាត់គំរូទី២ (MN = PQ ⇔ ធ្នូ MN = PQ) និងលំហាត់គំរូទី៣ (ធ្នូ AB > CD ⇒ AB > CD)
\newcommand{\figModelCIRCThree}{%
\begin{tikzpicture}[scale=0.38, every node/.style={font=\scriptsize, text=black}]
  % គំរូទី២ (ទំ.២០១)
  \begin{scope}[shift={(0,0)}]
    \coordinate (O) at (3.0,2.8);
    \def\R{2.35}
    \coordinate (M) at ($(O)+(130:\R)$);
    \coordinate (N) at ($(O)+(210:\R)$);
    \coordinate (P) at ($(O)+(50:\R)$);
    \coordinate (Q) at ($(O)+(-30:\R)$);
    \draw[black, line width=0.95pt] (O) circle (\R);
    \draw[black, line width=1.5pt] (M) arc (130:210:\R);
    \draw[black, line width=1.5pt] (Q) arc (-30:50:\R);
    \draw[black, line width=1.0pt] (M)--(N) (P)--(Q);
    \draw[black, line width=0.8pt] (O)--(M) (O)--(N) (O)--(P) (O)--(Q);
    \tickone{O}{M}\tickone{O}{N}\tickone{O}{P}\tickone{O}{Q}
    \ticktwo{M}{N}\ticktwo{P}{Q}
    \fill[black] (O) circle (2.2pt) (M) circle (2pt) (N) circle (2pt) (P) circle (2pt) (Q) circle (2pt);
    \node[below=2pt] at (O) {$O$};
    \node[above left=0.5pt] at (M) {$M$};
    \node[below left=0.5pt] at (N) {$N$};
    \node[above right=0.5pt] at (P) {$P$};
    \node[below right=0.5pt] at (Q) {$Q$};
    \node[align=center] at (3.0,-0.60) {គំរូទី២៖ $\tri{NOM} \cong \tri{POQ}$\\(ជ.ជ.ជ) $\Rightarrow \arc{MN} = \arc{PQ}$};
  \end{scope}
  % គំរូទី៣ (ទំ.២០១)
  \begin{scope}[shift={(9.3,0)}]
    \coordinate (O2) at (3.0,2.8);
    \def\R{2.35}
    \coordinate (A2) at ($(O2)+(115:\R)$);
    \coordinate (B2) at ($(O2)+(225:\R)$);
    \coordinate (C2) at ($(O2)+(35:\R)$);
    \coordinate (D2) at ($(O2)+(-20:\R)$);
    \draw[black, line width=0.95pt] (O2) circle (\R);
    \draw[black, line width=1.5pt] (A2) arc (115:225:\R);
    \draw[black, line width=1.5pt] (D2) arc (-20:35:\R);
    \draw[black, line width=1.0pt] (A2)--(B2) (C2)--(D2);
    \draw[black, line width=0.8pt] (O2)--(A2) (O2)--(B2) (O2)--(C2) (O2)--(D2);
    \tickone{O2}{A2}\tickone{O2}{B2}\tickone{O2}{C2}\tickone{O2}{D2}
    \fill[black] (O2) circle (2.2pt) (A2) circle (2pt) (B2) circle (2pt) (C2) circle (2pt) (D2) circle (2pt);
    \node[below=2pt] at (O2) {$O$};
    \node[above left=0.5pt] at (A2) {$A$};
    \node[left=1pt] at (B2) {$B$};
    \node[above right=0.5pt] at (C2) {$C$};
    \node[below right=0.5pt] at (D2) {$D$};
    \node[align=center] at (3.0,-0.60) {គំរូទី៣៖ $\angle AOB > \angle COD$\\$\Rightarrow AB > CD$ (ត្រចៀកទ្វារ)};
  \end{scope}
\end{tikzpicture}}

% ៤. រូប E5៖ លំហាត់ទី៣ (AB = CD ⇒ ធ្នូ AC = BD) និងលំហាត់ទី២ ទំ.២០៤
\newcommand{\figPracticeCIRCThree}{%
\begin{tikzpicture}[scale=0.38, every node/.style={font=\scriptsize, text=black}]
  % លំហាត់ទី៣ ទំ.២០៤
  \begin{scope}[shift={(0,0)}]
    \coordinate (O) at (3.0,2.8);
    \def\R{2.35}
    \coordinate (A) at ($(O)+(165:\R)$);
    \coordinate (B) at ($(O)+(105:\R)$);
    \coordinate (C) at ($(O)+(55:\R)$);
    \coordinate (D) at ($(O)+(-5:\R)$);
    \draw[black, line width=0.95pt] (O) circle (\R);
    \draw[black, line width=1.05pt] (A)--(B) (C)--(D);
    \draw[black, line width=1.5pt] (C) arc (55:105:\R);
    \ticktwo{A}{B}\ticktwo{C}{D}
    \fill[black] (O) circle (2.2pt) (A) circle (2pt) (B) circle (2pt) (C) circle (2pt) (D) circle (2pt);
    \node[below=1pt] at (O) {$O$};
    \node[left=0.5pt] at (A) {$A$};
    \node[above left=0.5pt] at (B) {$B$};
    \node[above right=0.5pt] at (C) {$C$};
    \node[right=0.5pt] at (D) {$D$};
    \node[align=center] at (3.0,-0.60) {លំហាត់ទី៣៖ $AB=CD$\\បូកធ្នូរួម $\arc{BC} \Rightarrow \arc{AC}=\arc{BD}$};
  \end{scope}
  % លំហាត់ទី២ ទំ.២០៤
  \begin{scope}[shift={(9.3,0)}]
    \coordinate (O2) at (3.0,2.8);
    \def\R{2.35}
    \coordinate (A2) at ($(O2)+(140:\R)$);
    \coordinate (B2) at ($(O2)+(40:\R)$);
    \coordinate (C2) at ($(O2)+(215:\R)$);
    \coordinate (D2) at ($(O2)+(-35:\R)$);
    \draw[black, line width=0.95pt] (O2) circle (\R);
    \draw[black, line width=0.95pt] (A2)--(D2) (B2)--(C2);
    \draw[black, line width=0.8pt, densely dashed] (A2)--(B2) (C2)--(D2);
    \fill[black] (O2) circle (2.2pt) (A2) circle (2pt) (B2) circle (2pt) (C2) circle (2pt) (D2) circle (2pt);
    \node[below=2pt] at (O2) {$O$};
    \node[above left=0.5pt] at (A2) {$A$};
    \node[above right=0.5pt] at (B2) {$B$};
    \node[below left=0.5pt] at (C2) {$C$};
    \node[below right=0.5pt] at (D2) {$D$};
    \node[align=center] at (3.0,-0.60) {លំហាត់ទី២ (កិច្ចការផ្ទះ)៖\\$\arc{AD} = \arc{CB} \Rightarrow AD = CB$};
  \end{scope}
\end{tikzpicture}}

\begin{document}
\fontsize{10pt}{15.0pt}\selectfont

% ==================== ក្បាលរដ្ឋបាលផ្លូវការ ====================
{\centering
  {\muol\fontsize{11pt}{16pt}\selectfont ព្រះរាជាណាចក្រកម្ពុជា}\\[2pt]
  {\muol\fontsize{11pt}{16pt}\selectfont ជាតិ សាសនា ព្រះមហាក្សត្រ}\\[-2pt]
  \khmflourish\par
}
\vspace{-4pt}
{\muol\fontsize{10pt}{15.2pt}\selectfont
មន្ទីរអប់រំ យុវជន និងកីឡា ខេត្តកំពង់ធំ\\
ការិយាល័យអប់រំ យុវជន និងកីឡា ស្រុកកំពង់ស្វាយ\\
អនុវិទ្យាល័យព្រះក្រសាំង\par}

\vspace{2pt}
{\centering{\muol\fontsize{12pt}{18pt}\selectfont កិច្ចតែងការបង្រៀន}\par}
\vspace{2pt}

{\fontsize{10pt}{15.0pt}\selectfont
\hspace*{3.7cm}%
\begin{tabular}{@{}l l l l@{}}
  - & មុខវិជ្ជា & : & គណិតវិទ្យា ថ្នាក់ទី៨ \\
  - & ចំណងជើងមេរៀន & : & មេរៀនទី១៥ រង្វង់ \\
  - & ចំណងជើងរង & : & ២.១ ធ្នូ និងអង្កត់ធ្នូ (ទំព័រ ២០០--២០១ និង ២០៤) \\
  - & រយៈពេល & : & ១ ម៉ោងសិក្សា (៥០ នាទី) \\
  - & បង្រៀនដោយ & : & ហុក ប៊ុនថែន \\
  - & កាលបរិច្ឆេទ & : & ថ្ងៃចន្ទ ទី២១ ខែធ្នូ ឆ្នាំ២០២៦ \\
\end{tabular}

\vspace{1pt}

% -------------------- I. វត្ថុបំណង --------------------
\textbf{I.} {\muol វត្ថុបំណង}\\[1pt]
\hspace*{0.55cm}%
\begin{tabular}{@{}l l L{14.2cm}@{}}
  \wcirc \hspace{0.25cm} វិជ្ជាសម្បទា & : & សិស្សប្រាប់ទ្រឹស្ដីបទអង្កត់ផ្ចិតជាអង្កត់ធ្នូធំជាងគេក្នុងរង្វង់ និងទំនាក់ទំនងទៅវិញទៅមករវាងធ្នូប៉ុនគ្នា/មិនប៉ុនគ្នា ជាមួយនឹងអង្កត់ធ្នូប៉ុនគ្នា/មិនប៉ុនគ្នាបានត្រឹមត្រូវ ។ \\
  \wcirc \hspace{0.25cm} បំណិនសម្បទា & : & សិស្សប្រើសញ្ញាសម្គាល់ស-ខ្មៅ និងធនាគារហេតុផល \cnum{1}--\cnum{5} ស្រាយបញ្ជាក់លក្ខណៈអង្កត់ផ្ចិត អង្កត់ធ្នូប៉ុនគ្នា និងអង្កត់ធ្នូមិនប៉ុនគ្នាលើសន្លឹកកិច្ចការបានត្រឹមត្រូវ ។ \\
  \wcirc \hspace{0.25cm} ចរិយាសម្បទា & : & សិស្សមានទម្លាប់ផ្ទៀងផ្ទាត់អំណះអំណាងតាមតក្កវិទ្យា (មិនសន្និដ្ឋានតាមការស្មានសមាមាត្រខុស) និងឱ្យតម្លៃលើការអនុវត្តធរណីមាត្រក្នុងសិប្បកម្មខ្មែរ ។ \\
\end{tabular}

\vspace{1pt}

% -------------------- II. សម្ភារៈឧបទេស --------------------
\textbf{II.} {\muol សម្ភារៈឧបទេស}\\[1pt]
\hspace*{0.55cm}\parbox[t]{18.2cm}{$\bullet$ \hspace{0.25cm} សម្រាប់គ្រូ៖ សៀវភៅសិក្សាគោលគណិតវិទ្យាថ្នាក់ទី៨ (ទំ.២០០--២០១ និង ២០៤) ក្ដារខៀន ដីស ដែកឈាន និងបន្ទាត់វែង (គ្មានស្លាយ \en{LCD}) ។\\
$\bullet$ \hspace{0.25cm} សម្រាប់សិស្ស៖ សន្លឹកកិច្ចការសិស្ស ស-ខ្មៅ (\en{Worksheet\_CIRC3} ២ ទំព័រពេញ) សៀវភៅគោល បន្ទាត់ក្រិត និងដែកឈាន ។}

\vspace{1pt}

% -------------------- III. វិធីសាស្រ្តបង្រៀន --------------------
\textbf{III.} {\muol វិធីសាស្រ្តបង្រៀន}\\[1pt]
\hspace*{0.45cm}%
\begin{tabular}{@{}l l L{14.85cm}@{}}
  $\bullet$ \hspace{0.25cm} គោលវិធី & : & សិស្សមជ្ឈមណ្ឌល (បង្រៀនដោយមិនប្រើស្លាយ \en{LCD} យកសន្លឹកកិច្ចការសម្រាយបញ្ជាក់ជាគោល) \\
  $\bullet$ \hspace{0.25cm} យុទ្ធវិធីបង្រៀន & : & គំរូ \en{5Es} រួមផ្សំសញ្ញាសម្គាល់ស-ខ្មៅ (\en{Monochrome Coding}) ធនាគារហេតុផល \cnum{1}--\cnum{5} និងឧទាហរណ៍ច្រានចោល \\
\end{tabular}
}

\vspace{1pt}

% -------------------- IV. ដំណើរការបង្រៀន --------------------
\textbf{IV.} {\muol ដំណើរការបង្រៀន}

\vspace{1pt}

\renewcommand{\arraystretch}{1.15}
\begin{longtable}{|L{5.50cm}|L{7.90cm}|L{4.80cm}|}
\hline
\rowcolor{headgray}
\multicolumn{1}{|c|}{\muol សកម្មភាពគ្រូ} &
\multicolumn{1}{c|}{\muol ខ្លឹមសារមេរៀន} &
\multicolumn{1}{c|}{\muol សកម្មភាពសិស្ស} \\
\hline
\endfirsthead

\hline
\rowcolor{headgray}
\multicolumn{1}{|c|}{\muol សកម្មភាពគ្រូ} &
\multicolumn{1}{c|}{\muol ខ្លឹមសារមេរៀន} &
\multicolumn{1}{c|}{\muol សកម្មភាពសិស្ស} \\
\hline
\endhead

\hline
\endfoot

\hline
\endlastfoot

% ==================== ជំហានទី១ (2min) ====================
\stepheader{ជំហានទី១ ៖ រដ្ឋបាលថ្នាក់ (២ នាទី)}
$\bullet$ ពិនិត្យអវត្តមាន សណ្ដាប់ធ្នាប់ និងការត្រៀមបន្ទាត់ក្រិត ដែកឈានរបស់សិស្ស\par\vspace{1pt}
$\bullet$ ចែកសន្លឹកកិច្ចការ \en{Worksheet\_CIRC3} (១ សន្លឹកសងខាង) ដល់សិស្សគ្រប់គ្នា
&
$\bullet$ ពិនិត្យវត្តមាន វិន័យ និងសម្ភារៈសិក្សា (សន្លឹកកិច្ចការ បន្ទាត់ក្រិត និងដែកឈាន)
&
$\circ$ ប្រធានថ្នាក់រាយការណ៍អវត្តមាន\par\vspace{1pt}
$\circ$ ទទួលសន្លឹកកិច្ចការ និងសរសេរឈ្មោះ \\[3pt]
\hline

% ==================== ជំហានទី២ (9min: ≤3 + 3 + E1 3) ====================
\stepheader{ជំហានទី២ ៖ កែកិច្ចការផ្ទះ រំលឹកមេរៀនចាស់ និងភ្ជាប់មេរៀនថ្មី (៩ នាទី)}
\timenote{(កែកិច្ចការផ្ទះ $\le 3\unit{min}$)}\par\vspace{1pt}
$\bullet$ ហៅសិស្ស ២ នាក់ឡើងបង្ហាញចម្លើយសង្ខេបនៃលំហាត់ទី១ និងលំហាត់ច្រាសទី២ ពីម៉ោងមុន
&
\textbf{កំណែកិច្ចការផ្ទះ៖}\par\vspace{1pt}
១. $C(O,18\unit{cm})$៖ ធ្នូធំ $\arc{ACB} = 360^\circ - 25^\circ = \mathbf{335^\circ}$ ; $P = 113.04\unit{cm}$ $\Rightarrow l_{AB} \approx \mathbf{7.85\unit{cm}},\; l_{CD} = \mathbf{37.68\unit{cm}}$ ។\par\vspace{1pt}
២. \textbf{ចំណោទច្រាស៖} $P = 94.2\unit{cm} \Rightarrow \angle MON = \frac{15.7 \times 360^\circ}{94.2} = \mathbf{60^\circ}$ ។
&
$\circ$ ឡើងកែលើក្ដារខៀន និងផ្ទៀងផ្ទាត់ចម្លើយក្នុងសៀវភៅ \\[3pt]

\timenote{(រំលឹកមេរៀនទាក់ទង $3\unit{min}$)}\par\vspace{1pt}
$\bullet$ សួររំលឹកវិសមភាពត្រីកោណ និងករណីប៉ុនគ្នាទី៣ នៃត្រីកោណ (ជ.ជ.ជ)
&
\textbf{រំលឹកវិសមភាពត្រីកោណ និងត្រីកោណប៉ុនគ្នា៖}\par\vspace{1pt}
$\bullet$ \textbf{វិសមភាពត្រីកោណ៖} ក្នុង $\tri{OCD}$ ផលបូកពីរជ្រុងតែងធំជាងជ្រុងទីបី៖ $OC + OD > CD$ ។
&
$\circ$ ឆ្លើយរំលឹកវិសមភាពត្រីកោណ $OC + OD > CD$ \\[3pt]

$\bullet$ សួរបន្តអំពីករណី ជ.ជ.ជ និងទ្រឹស្ដីបទត្រចៀកទ្វារ (Hinge Theorem)
&
$\bullet$ \textbf{ករណីទី៣ (ជ.ជ.ជ)៖} បើជ្រុងទាំងបីស្មើគ្នារៀងគ្នា នោះត្រីកោណទាំងពីរប៉ុនគ្នា ។\par\vspace{1pt}
$\bullet$ \textbf{ត្រចៀកទ្វារ៖} បើជ្រុងពីរស្មើគ្នារៀងគ្នា តែមុំចន្លោះធំជាង នោះជ្រុងឈមធំជាង ។
&
$\circ$ រំលឹកករណី ជ.ជ.ជ និងទ្រឹស្ដីបទត្រចៀកទ្វារ \\[3pt]

\fiveE{E1 - Engage}{3 min}\par\vspace{1pt}
$\bullet$ ឱ្យសិស្សពិនិត្យ \textbf{ផ្នែកទី១} លើសន្លឹកកិច្ចការ៖ ជាងត្បាញ \textbf{ចង្អេរឫស្សី} រាងរង្វង់ $C(O, 25\unit{cm})$ ចង់ដាក់របារឫស្សីត្រង់ ១ ដើមភ្ជាប់គែមចង្អេរឱ្យវែងបំផុត
&
\textbf{បេសកកម្ម E1៖ របារឫស្សីទ្រចង្អេរ ($r = 25\unit{cm}$)}\par\vspace{1pt}
$\bullet$ សំណួរបំផុស៖ (១) តើត្រូវដាក់របារត្រង់កាត់តាមចំណុចណា ទើបវែងជាងគេបំផុត? (២) បើធ្នូ $\arc{CF} = 2\arc{CD}$ តើអង្កត់ធ្នូ $CF = 2CD$ ឬ $CF < 2CD$? (ទុកផ្ទៀងនៅ E5)
&
$\circ$ សង្កេតរូបចង្អេរឫស្សី និងគូសធីកទស្សន៍ទាយទាំង ២ សំណួរលើសន្លឹកកិច្ចការ \\[3pt]

&
\cfig{\figEngageCIRCThree}
& \\[3pt]
\hline

% ==================== ជំហានទី៣ (17min: E2 11 + E3 6) ====================
\stepheader{ជំហានទី៣ ៖ រុករកលក្ខណៈធ្នូ និងអង្កត់ធ្នូតាមការវាស់មាត្រដ្ឋាន ១:១ (១៧ នាទី)}
\fiveE{E2 - Explore}{11 min}\par\vspace{1pt}
$\bullet$ ណែនាំដៃគូសិស្សប្រើបន្ទាត់ក្រិតវាស់ប្រវែងពិតលើរង្វង់ $C(O, 2\unit{cm})$ (មាត្រដ្ឋាន $1:1$) ក្នុង \textbf{ផ្នែកទី២}៖ អង្កត់ផ្ចិត $AB$ អង្កត់ធ្នូ $CD, EG$ ($60^\circ$) និង $CF$ ($120^\circ$)
&
\textbf{សកម្មភាពទី១៖ វាស់ប្រៀបធៀបលើរង្វង់ពិត $C(O, 2\unit{cm})$ (ទំ.២០០--២០១)}\par\vspace{1pt}
-- វាស់បាន៖ អង្កត់ផ្ចិត $AB = 4.0\unit{cm}$ ; អង្កត់ធ្នូ $CD = EG = 2.0\unit{cm}$ ; អង្កត់ធ្នូ $CF \approx 3.5\unit{cm}$ ។\par\vspace{1pt}
-- សង្កេតឃើញថា៖ អង្កត់ផ្ចិត $AB = 2r$ ធំជាងអង្កត់ធ្នូមិនកាត់ផ្ចិតទាំងអស់ ($AB > CF > CD$) ។
&
$\circ$ វាស់ប្រវែងពិតដោយបន្ទាត់ក្រិត និងកត់ត្រាចូលតារាងលើសន្លឹកកិច្ចការ \\[3pt]

$\bullet$ ឱ្យដៃគូប្រៀបធៀបធ្នូប៉ុនគ្នា ($\arc{CD} = \arc{EG}$) ធ្នូមិនប៉ុនគ្នា ($\arc{CF} > \arc{CD}$) និងពិនិត្យ \textbf{អន្ទាក់មិនសមាមាត្រ} ($\arc{CF} = 2\arc{CD}$)
&
\textbf{សកម្មភាពទី២៖ ទំនាក់ទំនងធ្នូ--អង្កត់ធ្នូ និងអន្ទាក់មិនសមាមាត្រ}\par\vspace{1pt}
-- ធ្នូប៉ុនគ្នា $\arc{CD} = \arc{EG} \; (60^\circ) \Leftrightarrow$ អង្កត់ធ្នូប៉ុនគ្នា $CD = EG \; (2.0\unit{cm})$ ។\par\vspace{1pt}
-- \textbf{អន្ទាក់ POE៖} ទោះបី $\arc{CF} = 2\arc{CD}$ ($120^\circ = 2\times 60^\circ$) ក៏ដោយ តែ $CF \approx 3.5\unit{cm} < 2CD = 4.0\unit{cm}$ (ព្រោះក្នុង $\tri{CDF}$ គេមាន $CF < CD + DF = 2CD$!) ។
&
$\circ$ រកឃើញថាអង្កត់ធ្នូមិនកើនឡើងសមាមាត្រពីរដងតាមធ្នូទេ ($CF < 2CD$) \\[3pt]

&
\cfig{\figExploreCIRCThree}
& \\[3pt]

\fiveE{E3 - Explain}{6 min}\par\vspace{1pt}
$\bullet$ សំយោគនិមិត្តសញ្ញា ($\arc{AB}$ ជាធ្នូ ឯ $AB$ ជាអង្កត់ធ្នូ) និងទ្រឹស្ដីបទទាំង ៣ តាមសៀវភៅគោល (ទំ.២០០--២០១)
&
\textbf{ទ្រឹស្ដីបទទាំង ៣ ស្ដីពីធ្នូ និងអង្កត់ធ្នូ (ទំ.២០១)៖}\par\vspace{1pt}
$\bullet$ \textbf{ទ្រឹស្ដីបទ ១៖} អង្កត់ផ្ចិតជាអង្កត់ធ្នូដែលធំជាងគេនៅក្នុងរង្វង់មួយ ។\par\vspace{1pt}
$\bullet$ \textbf{ទ្រឹស្ដីបទ ២៖} ក្នុងរង្វង់មួយ (ឬរង្វង់ពីរប៉ុនគ្នា) ធ្នូពីរប៉ុនគ្នាកំណត់ដោយអង្កត់ធ្នូពីរប៉ុនគ្នា និងច្រាសមកវិញ ($\arc{AB} = \arc{MN} \Leftrightarrow AB = MN$) ។
&
$\circ$ បំពេញទ្រឹស្ដីបទទី១ និងទី២ ក្នុង \textbf{ផ្នែកទី៣} នៃសន្លឹកកិច្ចការ \\[3pt]

$\bullet$ បន្តសំយោគទ្រឹស្ដីបទទី៣ និងបង្ហាញ \textbf{ធនាគារហេតុផល \cnum{1}--\cnum{5}} សម្រាប់ប្រើក្នុងសម្រាយបញ្ជាក់
&
$\bullet$ \textbf{ទ្រឹស្ដីបទ ៣៖} ក្នុងរង្វង់មួយ ធ្នូតូចពីរមិនប៉ុនគ្នា ធ្នូធំជាងកាត់ដោយអង្កត់ធ្នូធំជាង និងច្រាសមកវិញ ($\arc{AB} > \arc{CD} \Leftrightarrow AB > CD$) ។\par\vspace{1pt}
$\bullet$ \textbf{ធនាគារហេតុផល៖} \cnum{1} កាំរង្វង់តែមួយស្មើគ្នា $\;\vert\;$ \cnum{2} វិសមភាពត្រីកោណ $\;\vert\;$ \cnum{3} ត្រីកោណប៉ុនគ្នា (ជ.ជ.ជ) $\;\vert\;$ \cnum{4} មុំផ្ចិតស្មើគ្នា $\Leftrightarrow$ ធ្នូស្មើគ្នា $\;\vert\;$ \cnum{5} ទ្រឹស្ដីបទត្រចៀកទ្វារ ។
&
$\circ$ បំពេញទ្រឹស្ដីបទទី៣ និងកត់ត្រាធនាគារហេតុផល \cnum{1}--\cnum{5} \\[3pt]

$\bullet$ ឱ្យសិស្សសង្ខេបជាតារាង \textbf{សម្មតិកម្ម--សន្និដ្ឋាន} និងឆ្លើយពិត/មិនពិតលើសន្លឹកកិច្ចការ
&
{\renewcommand{\arraystretch}{1.12}\small
\begin{tabular}{|p{3.65cm}|p{3.65cm}|}
\hline
\cellcolor{headgray}\centering\textbf{សម្មតិកម្មក្នុង $C(O,r)$} &
\cellcolor{headgray}\centering\arraybackslash\textbf{សន្និដ្ឋាន / វិបាក} \\[2pt]
\hline
$AB$ ជាអង្កត់ផ្ចិត, $CD$ មិនកាត់ $O$ & $AB = 2r > CD$ (តាម \cnum{1}, \cnum{2}) \\[1.5pt]
\hline
ធ្នូតូច $\arc{MN} = \arc{PQ}$ & $MN = PQ$ និងច្រាសមកវិញ (\cnum{3}, \cnum{4}) \\[1.5pt]
\hline
ធ្នូតូច $\arc{AB} > \arc{CD}$ & $AB > CD$ និងច្រាសមកវិញ (\cnum{5}) \\[1.5pt]
\hline
\end{tabular}}
&
$\circ$ ផ្ទៀងផ្ទាត់តារាងសម្មតិកម្ម--សន្និដ្ឋាន និងឆ្លើយពិត/មិនពិត \\[3pt]
\hline

% ==================== ជំហានទី៤ (20min: E4 10 + E5 10) ====================
\stepheader{ជំហានទី៤ ៖ ពង្រឹងចំណេះដឹងតាមសម្រាយបញ្ជាក់លំហាត់គំរូ និងវាយតម្លៃ (២០ នាទី)}
\fiveE{E4 - Elaborate}{10 min}\par\vspace{1pt}
$\bullet$ ដឹកនាំសិស្សគូសសញ្ញាឆ្នូតស-ខ្មៅ (`|`, `||`) និងបំពេញលេខហេតុផលក្នុង \textbf{លំហាត់គំរូទី១} និង \textbf{លំហាត់គំរូទី២ (ទំ.២០១)} លើ \textbf{ផ្នែកទី៤}
&
\textbf{១. លំហាត់គំរូទី១ (ទំ.២០១ — អង្កត់ផ្ចិត $AB > CD$)៖}\par
-- ក្នុង $\tri{CDO}$ មាន $OC + OD > CD$ (តាម \cnum{2}) ។ ដោយ $OC = OD = r$ (តាម \cnum{1}) នោះ $OC + OD = 2r = AB \Rightarrow \mathbf{AB > CD}$ ។\par\vspace{1pt}
\textbf{២. លំហាត់គំរូទី២ (ទំ.២០១ — $MN = PQ \Rightarrow \arc{MN} = \arc{PQ}$)៖}\par
-- ក្នុង $\tri{NOM}$ និង $\tri{POQ}$ មាន $ON=OP=r,\; OM=OQ=r$ (\cnum{1}) និង $MN=PQ$ (សម្មតិកម្ម) $\Rightarrow \tri{NOM} \cong \tri{POQ}$ (\cnum{3}) $\Rightarrow \angle NOM = \angle POQ \Rightarrow \mathbf{\arc{MN} = \arc{PQ}}$ (\cnum{4}) ។
&
$\circ$ គូសឆ្នូតសម្គាល់កាំ និងបំពេញលេខហេតុផល \cnum{1}--\cnum{4} លើសន្លឹកកិច្ចការ \\[3pt]

$\bullet$ ឱ្យដៃគូបន្តបំពេញសម្រាយបញ្ជាក់ \textbf{លំហាត់គំរូទី៣ (ទំ.២០១)} ដោយប្រើទ្រឹស្ដីបទត្រចៀកទ្វារ \cnum{5}
&
\textbf{៣. លំហាត់គំរូទី៣ (ទំ.២០១ — ធ្នូតូច $\arc{AB} > \arc{CD} \Rightarrow AB > CD$)៖}\par
-- ដោយ $\arc{AB} > \arc{CD}$ គេបានមុំផ្ចិត $\angle AOB > \angle COD$ ។ ក្នុង $\tri{AOB}$ និង $\tri{COD}$ មាន $OA = OC = r,\; OB = OD = r$ (\cnum{1}) និង $\angle AOB > \angle COD \Rightarrow \mathbf{AB > CD}$ (តាម \cnum{5}) ។
&
$\circ$ បំពេញជំហានសម្រាយបញ្ជាក់គំរូទី៣ និងផ្ទៀងផ្ទាត់ជាមួយដៃគូ \\[3pt]

&
\cfig{\figModelCIRCThree}
& \\[3pt]

\fiveE{E5 - Evaluate}{10 min}\par\vspace{1pt}
$\bullet$ ឱ្យសិស្សស្រាយបញ្ជាក់ \textbf{លំហាត់ទី៣ (ទំ.២០៤)} ដោយខ្លួនឯងលើ \textbf{ផ្នែកទី៥} (គន្លឹះ៖ បូកធ្នូរួម $\arc{BC}$ ទាំងសងខាង)
&
\textbf{១. លំហាត់ទី៣ (ទំ.២០៤)៖} ក្នុងរង្វង់ $C(O,r)$ គេឱ្យអង្កត់ធ្នូ $AB = CD$ ។ បង្ហាញថា $\arc{AC} = \arc{BD}$ ។\par\vspace{1pt}
-- ដោយ $AB = CD \Rightarrow \arc{AB} = \arc{CD}$ (តាមទ្រឹស្ដីបទ ២) ។\par\vspace{1pt}
-- បូកធ្នូរួម $\arc{BC}$ អង្គទាំងពីរ៖ $\arc{AB} + \arc{BC} = \arc{CD} + \arc{BC} \Rightarrow \mathbf{\arc{AC} = \arc{BD}}$ ។
&
$\circ$ សរសេរសម្រាយបញ្ជាក់ដោយបូកធ្នូរួម $\arc{BC}$ លើសន្លឹកកិច្ចការ \\[3pt]

&
\cfig{\figPracticeCIRCThree}
& \\[3pt]

$\bullet$ ពិនិត្យចម្លើយបិទបេសកកម្មចង្អេរឫស្សី E1 ជាមួយសិស្សទូទាំងថ្នាក់
&
$\bullet$ \textbf{បិទបេសកកម្ម E1៖} (១) ត្រូវដាក់របារកាត់តាមផ្ចិត $O$ ជា \textbf{អង្កត់ផ្ចិត} $AB = 2r = 50\unit{cm}$ ទើបវែងបំផុត! (២) បើ $\arc{CF} = 2\arc{CD}$ នោះអង្កត់ធ្នូ $\mathbf{CF < 2CD}$!
&
$\circ$ ផ្ទៀងផ្ទាត់ការទស្សន៍ទាយដើមម៉ោង និងសរសេរចម្លើយបិទបេសកកម្ម \\[3pt]

$\bullet$ ធ្វើតេស្តរហ័ស ៦០ វិនាទី (\en{Exit Check}) ដោយឱ្យសិស្សលើកម្រាមដៃ ១--៣
&
$\bullet$ \textbf{Exit Check៖} ក្នុង $C(O, 6\unit{cm})$ អង្កត់ធ្នូធំបំផុត និងទំនាក់ទំនងពេល $\arc{PQ} = 2\arc{RS}$ គឺ៖ \textbf{A. $12\unit{cm}$ និង $PQ < 2RS$ (ត្រូវ)} $\;\vert\;$ B. $12\unit{cm}$ និង $PQ = 2RS$ $\;\vert\;$ C. $6\unit{cm}$ ។
&
$\circ$ លើកម្រាមដៃទី១ ជ្រើសចម្លើយត្រឹមត្រូវ \textbf{A} \\[3pt]
\hline

% ==================== ជំហានទី៥ (≤2min) ====================
\stepheader{ជំហានទី៥ ៖ កិច្ចការផ្ទះ និងបណ្ដាំផ្ញើ ($\le$ ២ នាទី)}
\timenote{(ចំណាយពេល $\le 2\unit{min}$)}\par\vspace{1pt}
$\bullet$ ដាក់កិច្ចការផ្ទះ លំហាត់ទី២ និងទី៥ (ទំ.២០៤)\par\vspace{1pt}
$\bullet$ ផ្ដាំផ្ញើឱ្យអានមុន ចំណុច \textbf{២.២ លក្ខណៈឆ្លុះនៃរង្វង់} (ទំ.២០២--២០៣) និងយកបន្ទាត់កែងមកម៉ោងក្រោយ
&
\textbf{កិច្ចការផ្ទះ (សៀវភៅសិក្សាគោល ទំព័រ ២០៤)៖}\par\vspace{1pt}
១. \textbf{លំហាត់ទី ២៖} ក្នុងរង្វង់ $C(O,r)$ តាមរូបគេមានធ្នូ $\arc{AD} = \arc{CB}$ ។ ចូរបង្ហាញថាអង្កត់ធ្នូ $AD = CB$ (ព្រមទាំងពិនិត្យករណីបើ $\arc{AC} = \arc{BD}$ បូកធ្នូរួម $\arc{CD}$) ។\par\vspace{1pt}
២. \textbf{លំហាត់ទី ៥៖} ក្នុងរង្វង់ $C(O,r)$ គេឱ្យអង្កត់ធ្នូ $AB = CD$ ។ ចូរបង្ហាញថាធ្នូ $\arc{BD} = \arc{AC}$ ។
&
$\circ$ កត់ត្រាកិច្ចការផ្ទះ និងត្រៀមអានចំណុច ២.២ មុន\par\vspace{1pt}
$\circ$ ទទួលបណ្ដាំផ្ញើ និងជម្រាបលាគ្រូ \\[3pt]
\end{longtable}

% -------------------- កាលបរិច្ឆេទ និងហត្ថលេខា --------------------
\vspace{8pt}
\hfill
\begin{minipage}{9.8cm}
  {\centering
    ថ្ងៃចន្ទ ១២កើត ខែមិគសិរ ឆ្នាំមមី អដ្ឋស័ក ពុទ្ធសករាជ ២៥៧០\\[2pt]
    ព្រះក្រសាំង, ថ្ងៃទី២១ ខែធ្នូ ឆ្នាំ២០២៦\\[3pt]
    {\muol គ្រូបង្រៀន}\par}
  \vspace{26pt}
  \hfill{\muol ហុក ប៊ុនថែន}\hspace*{1.4cm}
\end{minipage}

\end{document}
